\begin{textsample}{POS Dim 4 – vem\_ed\_12 – Score 3.00 – vem\_ed\_12/t050.txt}  \label{ex:f4_pos_141}
QUEBRA-GELO Prospectar novas colaborações, garantir uma média semestral de 25 Projetos Colaborativos Internacionais (PCIs / Cesu) e manter a “boa qualidade dos PCIs relatada pelos participantes por meio de Pesquisas de Percepção aplicadas \textbf{junto} a alunos e professores participantes ao final dos projetos”.
Eis algumas das metas do Programa de Línguas e Projetos Internacionais no Plano de Metas 2022 da Unidade do Ensino Superior de Graduação (Cesu).
Para avaliar a qualidade dos projetos, a equipe dos PCIs / Cesu desenvolveu um instrumento de pesquisa para verificar a percepção de professores e estudantes envolvidos nos PCIs, ao fim de cada semestre letivo.
Os questionários são aplicados na plataforma Microsoft Forms desde o segundo semestre de 2020 e a boa avaliação \textbf{vem} se mantendo ao longo do tempo.
Nas páginas 2 e 3 desta edição, você pode conferir uma síntese dos principais resultados colhidos no segundo semestre de 2021.
Por exemplo: 96 \% dos alunos \textbf{recomendariam} aos colegas realizar um PCI e 93 \% acreditam que participar do projeto melhora o desempenho acadêmico e as chances no mercado de trabalho.
Todos professores que responderam à pesquisa \textbf{recomendam} os PCIs para os colegas e concordam que a iniciativa possibilita realizar novas atividades acadêmicas, além de melhorar a interação com seus alunos.
A equipe dos PCIs \textbf{vem} trabalhando para fortalecer parcerias com Portugal.
É o caso da Universidade de Aveiro, que desde o segundo semestre de 2021 desenvolve PCIs com três Fatecs – Guaratinguetá, Lins e Ipiranga –, nas áreas de comunicação empresarial, turismo e marketing.
Essas colaborações estão descritas na página 4.
Boa leitura!

% matched lemmas: junto, recomendar, vir
\end{textsample}
